\section{Local sector factors and their complete supports}
\label{app:rg-geometry}

This section makes the geometric part of the block integration explicit.
The energy rows, the truncated free operators, and their estimates are
those of Section~\ref{rg:gaussian}.  We prove that the factors can be
prepared without changing a sector integral, that the small logarithm
persists where its derivatives are used, and that separated factors
integrate independently.  Positivity and exponential localization of the
free operators remain inputs from Lemma~\ref{lem:free-stack}; the present
argument does not prove those analytic estimates.

\subsection{One metric and one assignment of each factor}
\label{app:factor-ownership}

Fix $L$ before making the following convention.  Embed each coarse
position at its block anchor, and measure all distances in the same
fine-lattice metric.  The finite local formulas use a fixed number of
primitive kernels.  Choose their truncation radii, and the inverse-window
radius $R$, small enough that the sum of the radii in any such local
composition is at most $M/100$.  Include block diameters, reference-point
displacements, and kinetic-mask radii in this sum.  This is possible by
reducing the kernel cutoff by a fixed factor depending on $L$, and then
increasing $H$: the block diameter is fixed and the mask radius is
$O_L(\log H_j)=o(M)$.  The omitted-kernel bounds remain
$C_Le^{-c_LM}$ with a possibly smaller $c_L>0$.  Long convolution tails
and determinant or inverse chains are not local compositions for this
convention: every successive endpoint and window is retained in their
records.

For a set $P$ of seed anchors, let $N_r(P)$ denote its closed
$r$-neighborhood.  Seeds have their edge, observation, or output-edge
types attached.  Join two seeds when their $50M$-neighborhoods intersect.
Write $P_c$ for the seed set of a resulting component.  Distinct sets
$N_{50M}(P_c)$ are disjoint.  Consequently each of the assignments below
has at most one owner:
\begin{center}
\begin{tabular}{ll}
\toprule
Object & Owner $c$, if its anchor lies in the indicated set\\
\midrule
Frozen spin & $N_{3M}(P_c)$\\
Numerical energy row & $N_{5M}(P_c)$\\
Stationary square & $N_{8M}(P_c)$\\
Edge or observation default & $N_{14M}(P_c)$\\
Leading determinant change & $N_{3M+R}(P_c)$\\
\bottomrule
\end{tabular}
\end{center}
The last assignment is applicable because an inverse window differs
from the empty geometry only when it sees a frozen site.  Objects
without an owner remain separate factors.  Every primary profile belongs
to its own component.  Every frozen-site normalization reciprocal belongs
to that site's owner.  This assigns each additive row and each
multiplicative test once; in particular, the $8M$ and $14M$ collections
are not additional copies of the $5M$ collection.

The maximality of $P_c$ is implemented by compatibility of the
$50M$ footprints in the primary-pattern sum.  It is not an additional
predicate in the core factor.  Testing maximality directly could read a
seed $100M$ away.  The pattern is fixed throughout Gaussian integration;
it is never recomputed from a Gaussian value.

\subsection{Profiles and the principal logarithm}

Choose once and for all $\vartheta\in C^\infty(\mathbb R,[0,1])$
equal to one on $(-\infty,0]$ and zero on $[1,\infty)$.  Define
\[
 G_e(q)=\vartheta\left(
 \frac{|d_eq|^2-(t/4)^2}{(t/2)^2-(t/4)^2}\right),\qquad
 G_Y(q,V)=\vartheta\left(
 \frac{|V_Y-\widehat m_Y(q)|^2-(Wt)^2}{3(Wt)^2}\right).
\]
The tag in $\widehat m_Y$ is retained wherever its fallback branch is
used.  The factor for a selected primary is $1-G$; otherwise it is $G$.
These give exactly the sector decomposition by $G+(1-G)=1$.  For every
fixed derivative order, a derivative of $G_e$ has support in
$\{|d_eq|\le t/2\}$, and a derivative of $1-G_e$ has support in
$\{|d_eq|\ge t/4\}$.  The corresponding observation bounds are $2Wt$
and $Wt$.  Here and below derivatives are taken only on a smooth branch;
measurable physical-spin dependencies are held fixed by the transport
identity~\eqref{rg:transport}.

\begin{lemma}[Alignment and its derivative supports]
\label{app:principal-log}
Suppose that all internal edge defaults of an $L$-block and its
observation default are retained in a product.  On the support of this
product, or of any product-rule term of a fixed-order derivative of its
smooth profiles, every constituent spin satisfies
\[
 |q_xV_Y^{-1}-1|\le K_Lt.
\]
If $q_x=\exp(\xi_x)V_Y$ and $|\xi_x|<\pi$, then
$\xi_x$ is the principal logarithm and
$|\xi_x|\le 2K_Lt$ for sufficiently large $H$.
\end{lemma}
\begin{proof}
A coordinate path in the block has length at most $2L$.  The retained
edge bounds imply $|q_x-q_z|\le Lt$ for all constituent spins.  As the
block weights are nonnegative and sum to one,
$|q_x-m_Y|\le Lt$ and $|m_Y|\ge1-Lt$.  For $Lt<1/4$, the normalized
mean branch is therefore used and
$|q_x-m_Y/|m_Y||\le2Lt$.  The retained observation bound gives
$|q_x-V_Y|\le(2L+2W)t$.  Product-rule derivatives preserve the closed
support bounds of every profile, so the argument is unchanged there.
For a unit quaternion,
$|\exp\xi-1|=2\sin(|\xi|/2)$ when $0\le|\xi|<\pi$.
The asserted logarithm and its bound follow from the inverse sine.
\end{proof}

Every exterior spin is farther than $3M$ from the seeds.  Its block
therefore contains no edge or observation primary; the locality
convention makes its diameter less than $M/100$.  Lemma~\ref{app:principal-log}
applies on the full sector.  It also applies locally to every spin of a
selected assigned row: that row is farther than $2.5M$ from every seed,
and its block tests lie inside the retained $14M$ neighborhood.

Here is a completely smooth extension convention.  Put $a_t=2K_Lt$,
enlarging $K_L$ if the finitely many row reference points require it.
On $S^3$, use a smooth function of the geodesic angle that equals one
through $a_t$ and zero from $2a_t$, multiply the principal logarithm
by it, and extend the product by zero.  This defines a global smooth
frozen-angle function, since $2a_t<\pi/2$.  It equals the actual frozen
angle on every selected-row support just described.

On an exterior coordinate use a protector equal to one through
$|\xi|=\pi/2$ and zero from $3\pi/4$.  Also, when that coordinate is
read by a core factor or another protected hard factor, attach an
alignment profile equal to one through $a_t$ and zero from $2a_t$.
Both profiles are smooth functions of $|\xi|^2$ and remain attached.
For a compulsory core or mandatory old covering weight they multiply
the factor itself.  For an optional hard exponential $e^V$, first open
it and prepare instead
\[
 1+\chi_{\rm att}(e^V-1),
\]
where $\chi_{\rm att}$ is the product of its attached profiles.
Thus the profiles multiply its selected letter, including an old
nonremote masked term or a hard kinetic tail.  This equals $e^V$ on
the full sector and retains the small bound on $e^V-1$.
Opening $\chi_{\rm att}e^V$ instead would introduce an unpriced
order-one failure outside its plateau, and is not the prescription.
The alignment profile is the explicit larger smooth
plateau allowed in the preparation; it imposes no hard moving-angle
predicate.  It is one on the full sector and on the local supports needed
for the reserve.  Its usefulness at the outer boundary of the retained
defaults is that a spin read by a boundary default need not have every
test of its block retained.  Every such read is nevertheless in a small
chart on the support of the attached factor and all its derivatives.

The separate global chart profile is one through $\pi/4$ and zero
from $\pi/3$.  Multiplication by this profile before extending the
angular integral to $\mathbb R^3$ changes no sector integral, by
Lemma~\ref{app:principal-log}.  It removes the other exponential
preimages.  Its subsequent expansion as one plus a failure factor does
not detach the protectors or alignment profiles just specified.

\subsection{An explicit core factor and its one-time energy charges}

Let $J_c$ be the numerical rows assigned to $c$, $J_c^{\rm sel}$
the selected $\mathcal F,\mathcal B$ rows among them, and $S_c$ the
stationary-square rows assigned to $c$.  Write $L_i$ for the appropriate
row among $\mathcal F,\mathcal B,\mathcal C,\mathcal U,\mathcal N,
\mathcal Z$.  Put
\[
 K_c=\sum_{i\in J_c}L_i
      -\frac12\sum_{i\in J_c^{\rm sel}}|D_i\xi+F_i|^2,
 \qquad
 H_c=\frac12\sum_{i\in S_c}|r_{s,i}|^2.
\]
All rows in the first sum are included with their original sign.
In particular no second copy of a direct energy reserve is introduced
for a mask witness or a mean-error witness.

Let $\mathcal P_c$ be the product of the primary profiles of $c$ and
the defaults assigned to $c$.  Let $\mathcal I_c$ retain the local
inventory checks: a marked true bad input edge is supplied once by its
old covering label, and every unmarked selected input primary has the
factor $\mathbf1_{r_e\le t}$.  Retain each queried output status,
including an absent status, in the formula's data.  For an output-empty
extension replace these output predicates by the analytic no-flag
formula on the stated output graph domain, as in Section~\ref{rg:gaussian}.
This is an extension, not a differentiation across the jump of a status
indicator.

The leading determinant normalization can be written without an
unspecified local compensation.  The quadratic acts identically in the
three tangent colors.  For an exterior site $x$ define
\[
 h_x=-\frac32\int_0^\infty
       \left((1+u)^{-1}-(A_{\Omega_x}+u)^{-1}_{xx}\right)\,du,
\]
and let $h_\circ(x)$ be its empty-pattern value at the same fine-site
phase modulo the blocking lattice.  This phase must be retained:
$Q^*Q$ need not be invariant under one-step fine translations.
Let $F_c$ now denote
the set of frozen sites of $c$ (distinct from the affine vector $F_i$).
With $\kappa_g=(2\pi g^2)^{3/2}J_{\exp}(0)$, define
\[
 N_c=\kappa_g^{-|F_c|}
 \exp\left\{\sum_{\substack{x\text{ exterior}\\
                    x\in N_{3M+R}(P_c)}}(h_x-h_\circ(x))
                    -\sum_{x\in F_c}h_\circ(x)\right\}.
\]
The sum includes zero differences harmlessly.  Uniform upper and lower
spectral bounds on $A_{\Omega_x}$ give $|h_x|\le C_L$; at infinity
the integrand is $O_L((1+u)^{-2})$.  Thus
$|\log N_c|\le C_LM^2(1+|\log g|)|P_c|$.
This is exactly the correction obtained by extracting
$\kappa_g^{|\Lambda|}\exp\{\sum_{x\in\Lambda}h_\circ(x)\}$
while retaining Haar integration at frozen sites.  The sum is a fixed
scalar per complete $L$-block, so it has the required volume form.
The residual determinant chains stay
separate, and the observation denominator is extracted on every block.

Writing $\mathcal A_c$ for the product of the attached protectors and
alignment profiles, one possible compulsory core factor is therefore
\begin{equation}
 B_c=\mathcal P_c\mathcal I_c\mathcal A_c\,N_c
                 \exp\{-b(K_c+H_c)\}.
 \label{app:explicit-core}
\end{equation}
The remaining factors are exactly the rows, stationary squares,
defaults, Jacobian ratios, and determinant or inverse chains without an
owner, and the separately supported old terms.  Subtracting every
selected affine square once and multiplying back its Gaussian integral
is an identity.  The assignments, together with the plateau equalities
proved above, consequently preserve the full sector integral exactly.

For clarity we give the charging argument for the reserve in this
notation.  The quantitative inputs are the stack and ledger bounds
\eqref{rg:stack-sizes}--\eqref{rg:ledger-errors}, with their weighted
row and column sums.  An input seed with $t/4\le r\le t$ has a
direct square at least $ct^2$; for $r>t$ its direct row is at least
$ctr$ once $L$ is sufficiently large.  A noise seed has
$|Y-w|^2/2\ge ct^2$ once $W$ is large.  An output seed has the positive
direct $\mathcal C$ reserve displayed in Section~\ref{rg:gaussian}.
The second-stack $\mathcal C$ contribution cancels and $\mathcal N$
is nonnegative.

Only a failed input kinetic mask can make a second-stack
$\mathcal F$ contribution negative.  Charge that contribution to all
its bad-edge witnesses with the following upper bound, so no choice of
witness or charge multiplicity is hidden.  A witness $f$ can occur in
at most $C(1+\log(r_f/t))^2\le Cr_f/t$ mask tests.  Splitting the
convolution into chords at most $t$ and chords greater than $t$, and
using its column sum for the latter, bounds the total negative
contribution by
\[
 \frac{Ct}{L}\sum_{f:r_f>t}r_f.
\]
Every witness of an assigned row belongs to the same primary component:
the mask radius is smaller than $M/100$ and that witness is itself an
input seed.  Its direct row is therefore assigned to this component.
Choose $L$ so the displayed loss consumes less than a fixed fraction
of those direct reserves.

For the mean term, divide the observation blocks read by assigned rows
into blocks all of whose chords are at most $T$ and their complement.
The former cost at most $C_LtT^2$ per block.  In the latter, use
$|m-\widehat m|\le (C/L)\sum_{e\subset Y}r_e$.  Chords greater
than $t$ cost another $(Ct/L)\sum_{r_f>t}r_f$.  Chords at most $t$
cost at most $C_Lt^2$ per such block.  Each such block has a witness
$r_f>T$, so their number is bounded, with fixed multiplicity, by
$T^{-1}\sum_{r_f>T}r_f$.  This last cost is absorbed by the unused
$ctr_f$ reserve because $t/T\to0$.  This uses the same reserves with
specified fractional budgets, rather than paying for a witness twice.

Selected assigned rows are farther than $2.5M$ from every seed.  Their
masks are enabled, their means use the smooth branch, and no row in
this collection supplies a seed or failed-mask reserve.  Their small
chart bounds and Taylor's formula give
\[
 |R_i-R_i^{\rm aff}|\le C_LM^Dt^2,
 \qquad |R_i|+|R_i^{\rm aff}|\le C_LM^Dt.
\]
Thus removal of their squares costs $C_LM^Dt^3$ per row, independently
of the preceding charges.  There are at most $C_LM^2|P_c|$ rows,
blocks, or local tests in any fixed neighborhood of $P_c$.  This
follows by covering that neighborhood by one fixed-radius ball about
each seed; overlap reduces, rather than increases, the count.  The
number $D$ of polynomial losses is fixed by the finite local formulas
and the fixed derivative order, not by the geometry of $P_c$.
Numerical truncation and the remaining small-chord mean terms therefore
give
\[
 K_c\ge c t^2|P_c|
 -C_LM^D|P_c|\{t^3+tT^2+te^{-c_LM}\}.
\]
Increasing $H$ absorbs the last line as in the stated scale choices.
Since $H_c\ge0$, \eqref{app:explicit-core} then gives
\[
 |B_c|\le
 \exp\{-c_Lp^2|P_c|+C_LM^D(1+|\log g|)|P_c|\}.
\]
The same calculation holds on the supports of differentiated retained
profiles: primary lower bounds and default upper bounds persist,
and the selected angles remain the principal angles by
Lemma~\ref{app:principal-log}.  Alignment profiles additionally control
outer-boundary reads.  Product-rule differentiation does not discard an
undifferentiated test; it replaces it by a derivative with support in
the same closed support.  Hard physical-spin tests are transported,
not differentiated.  The no-output-flag extension uses the positive-mask,
unclipped formula through $4t'$; it changes only the fixed constants in
the chart and Taylor estimates.

\subsection{The read sets of the prepared factors}
\label{app:read-sets}

The following conventions specify complete read sets.  A raw position
means every endpoint of a row stencil, graph, or recorded path, rather
than only its anchor.  Add the mean window and output reference for each
exterior spin read.  Add both centered inverse windows for each covariance
entry used.  Add the $3M$ seed-status neighborhood needed to decide
membership in every such window.  A status read includes both possible
answers and is retained when the factor is considered alone.

\begin{center}\small
\begin{tabular}{>{\raggedright\arraybackslash}p{.24\textwidth}p{.68\textwidth}}
\toprule
Factor & Raw arguments and additional queries\\
\midrule
$B_c$ & All assigned rows, $8M$ stationary rows, $14M$ defaults,
frozen sites, local inventories, tags, and every queried input/output
mask or status; the determinant windows in $N_c$.\\[3pt]
Unassigned $\mathcal F,\mathcal B,\mathcal U$ & All row endpoints,
mean-block spins and tags, output references, clipping/mask statuses
or the specified analytic no-flag branch; cutoff arguments.\\[3pt]
Stationary, $\mathcal N,\mathcal Z$ & All numerical stencil endpoints;
the affine-vector arguments in $D\mu+F$; their seed-selection queries.\\[3pt]
Old canonical term or error & Its entire old graph, anchored joining
paths, compensation terms, masks, and the absence-of-seeds query in
$N_{14M}$ of the projected graph that defines ``remote.''\\[3pt]
Kinetic tail & Both complete paths, all endpoints, and each input/output
mask query.  A long path is never replaced by its anchor.\\[3pt]
Ratio or determinant chain & Every successive matrix endpoint and
window, including its hole tests; the affine-vector arguments for a
linear chain.  Keep the Gaussian endpoint set separately.\\[3pt]
Jacobian or chart factor & Its exterior coordinate and output reference;
its mean window and hole tests.\\[3pt]
Default or failure & Both edge endpoints or the whole observation block;
the tag if used, smooth mean-branch profile, and its seed/default status.\\[3pt]
Old covering weight & Its original complete support and inventory,
every physical-spin argument, and all attached profiles.  It retains its
own support even when it meets a core.\\
\bottomrule
\end{tabular}
\end{center}

A noncore factor retains $N_{20M}$ of every raw position and its
joining record.  The $14M$ remote query, $3M$ hole query, and finite
local composition enlargement fit in $18M$ under the chosen convention.
For a core all raw local tests lie in $N_{14M}(P_c)$; their enlargement
fits in $N_{18M}(P_c)\subset N_{50M}(P_c)$.  Old nonremote factors,
long tails, and long chains keep their own complete supports, and meet
a core or other factor by support intersection.  They are not absorbed
into a fixed-radius core while forgetting their distant endpoints.

\subsection{Exact local marginals}

Let $E(P)=\Lambda\setminus N_{3M}(P)$ and write $A_0$ for the
empty-pattern quadratic.  Every row incident on an exterior column is
selected: its anchor is farther than $3M-M/100>2.5M$ from the seeds.
Consequently
\[
 A=A_0[E(P),E(P)].
\]
There is no additional selected-row dependence in this principal matrix.
For $x\in E(P)$ let $\Omega_x=E(P)\cap B_R(x)$ and write
\[
 \Pi_P(y,x)=\mathbf1_{y\in\Omega_x}
       \bigl(A_0[\Omega_x,\Omega_x]^{-1}\bigr)_{yx},
 \qquad C_P=\tfrac12(\Pi_P+\Pi_P^*).
\]
The positivity of $C_P$ is supplied by the window estimate in
Section~\ref{rg:gaussian}.

\begin{lemma}[Local marginals and factorization]
\label{app:local-marginals}
The covariance $C_P(x,y)$ vanishes for $d(x,y)>R$.  For
$I\subset E(P)$ its marginal matrix $C_P[I,I]$ depends only on
$P\cap N_{R+3M}(I)$.  Its numerical mean on $I$ also depends only
on the mean windows, affine-row field arguments, and hole queries
specified above.

Suppose finitely many prepared factors have disjoint complete supports.
Their Gaussian coordinate sets have distance greater than $R$, their
frozen-spin and tag argument sets are disjoint, and their integrals
factor exactly.  Each individual integral is unchanged on deleting
outside primary components whose complete supports are disjoint from
its own, with all recorded external fields and status answers fixed.
\end{lemma}
\begin{proof}
A column of $\Pi_P$ vanishes outside its radius-$R$ window;
transposition preserves this assertion for an entry's two endpoints.
To decide $\Omega_x$ requires only seed positions within distance
$R+3M$ of $x$.  Once that set is known, the principal matrix and its
inverse are fixed.  Including the windows centered at both endpoints
therefore proves the assertion for $C_P[I,I]$.

The formula
\[
 \mu_x=-\sum_{y\in\Omega_x}
      \bigl(A_0[\Omega_x,\Omega_x]^{-1}\bigr)_{yx}(D^*F)_y
\]
has exactly the mean reads specified in the table.  Notice that its
inverse is a window inverse, not the full inverse $A^{-1}$.

The complete-support margins include an $R$-neighborhood of each
Gaussian read.  Disjoint complete supports consequently imply zero
cross covariance.  The characteristic function of the Gaussian
marginal on their union factors into the characteristic functions of
its blocks.  The frozen Haar measures and tag measures are independent
product measures, also independent of the Gaussian, so the complete
integrals factor.  All deterministic formulas and all marginal entries
remain unchanged when outside components are deleted, by the read-set
statements just proved.  These are Gaussian marginals with covariance
$C_P[I,I]$, not conditional distributions or principal restrictions of
the precision matrix.  Empty Gaussian or frozen argument sets are
allowed and have the unit product measure.
\end{proof}

Finally, at a fixed total derivative order only a fixed number of factors
are differentiated.  The number of placements is polynomial in the
number of read positions, hence in the total support load times a fixed
power of $M$.  For completeness the vector field used in protected
transport is, on each protected coordinate,
\[
 W_x=-g^{-1}(d\exp_{\xi_x})^{-1}
             \bigl((\partial_vq_x)T_x^{-1}\bigr),
 \qquad \xi_x=\mu_x+gZ_x.
\]
Here $\partial_v$ holds $Z$ fixed, and the inverse maps the tangent
space at $\exp\xi_x$ to the imaginary quaternions.  Extend this field
smoothly from the compact protector chart and set its other components
to zero.  All fixed-order derivatives of this inverse differential
are bounded on that compact chart.  The field derivatives of $\mu$
come from the displayed window formula: its inverse matrix is independent
of the output field at fixed primary statuses, and the derivatives of
its affine arguments use only the smooth finite-row and frozen-angle
formulas.  Their cutoffs cost fixed powers of $t^{-1}$, paired with
normalized directions of size $t$ times the prescribed distance weights.
Window and free-history derivative or difference bounds supply the
corresponding parameter estimates.  These facts, and the stated
prediction bounds for remote rows, give fixed polynomial bounds for
$W$, its divergence, and their iterated derivatives in $g^{-1},M,p$
and the number of reads.

At fixed geometry $C_P$ has no output-field dependence.  For a parameter
or history comparison interpolate its two positive marginal matrices
linearly.  Uniform ellipticity is preserved, and the covariance score is
\[
 \frac12\left(Z[I]^*C[I,I]^{-1}\dot C[I,I]
                   C[I,I]^{-1}Z[I]
             -\operatorname{Tr}(C[I,I]^{-1}\dot C[I,I])\right).
\]
The window difference estimate bounds $\dot C$ with its own discrepancy
factor.  The trace costs at most a constant times $|I|$, and the other
term is a Gaussian polynomial of fixed degree.  Thus the derivative
scores have the polynomial load dependence needed for the joint
estimate, with the discrepancy retained.  Every hard nonremote factor still meets a retained
core, an old covering weight carries a primary inventory edge, and a
voluntarily hard tail retains its exceptional chain estimate.  This
identifies the reserve used by the joint-factor estimate; it does not
replace that estimate by a claim of independence for overlapping
supports.
